% Chained immediate transfer --- plan \S4.8 (Phase 2).
% Sources: N=2 Immediate Transfer content/N_eq_2.tex (historical draft);
% results proved and verified 2026-08-07 by
% N=2 Immediate Transfer content/verify_chained_transfer.py (28/28 PASS).
% Everything displayed here is verified; the suite lives alongside this
% chapter's source folder.

\section{Chained immediate transfer}
\label{sec:chained}

\subsection{Setup}

We now couple several cells together in the most austere way possible. Take
$M$ host cells, all empty save one, which holds a single particle. Each cell
runs the $\mu=0$ birth--catastrophe process of the previous sections (birth
$\beta n$, catastrophe $\delta n$); on rupture, however, the released load
does not enter an extracellular pool. It is transferred immediately and
wholesale into a fresh cell, where a new BDC process begins; and so on,
until every cell has ruptured. There is no extracellular phase, no
clearance, and no re-infection --- only a chain of bursts.

The motivation is twofold. First, founding censuses can be remarkably small:
in plague, say, colony-forming units are counted in the hundreds against a
vast population of macrophages, so a serial, one-cell-at-a-time picture is
not absurd as a caricature of early pathogenesis. Second, the model is a
tractable laboratory for multi-cell burst statistics --- the sizes and times
of successive ruptures --- which interpolate between the single-cell theory
above and the population models of Part II. Since $\mu=0$, internal
extinction is impossible and all $M$ ruptures occur almost surely.

Let $r_i$ be the load at the $i$th rupture (the $i$th \emph{rupture size}),
$t_i$ its time, and $T_i=t_i-t_{i-1}$ the $i$th inter-rupture interval
(with $t_0=0$). Write
\begin{equation}
s=\frac{\delta}{\beta+\delta},
\qquad
\rho=\frac{\beta}{\beta+\delta}=1-s.
\end{equation}

\subsection{The key observation: event types ignore the load}
\label{sec:key-observation}

When the load is $n$, the next event is a birth with probability
$\beta n/(\beta n+\delta n)=\rho$ and a catastrophe with probability $s$
--- independent of $n$. Only the holding time depends on the load: it is
$\mathrm{Exp}\lrb{(\beta+\delta)n}$. Consequently the sequence of event
\emph{types} is a sequence of independent coin flips, and the number of
births $G$ that occur before the catastrophe in a cell founded with any load
is
\begin{equation}
G\sim\mathrm{Geom}_0(s),
\qquad
\pr{G=g}=s\,\rho^g,\qquad g=0,1,2,\ldots
\end{equation}
(the number of failures before the first success). Each cell in the chain
contributes an \emph{independent} copy $G_i$ of this variable --- the Markov
process restarts afresh at every transfer --- and the rupture size of cell
$i$ is its founding load plus its own births. Since the first cell is
founded by a single particle,
\begin{equation}
r_k = 1 + G_1 + \cdots + G_k,
\label{eq:rkdecomp}
\end{equation}
with $G_1,\ldots,G_k$ i.i.d.\ $\mathrm{Geom}_0(s)$. Everything about the
rupture sizes follows from this decomposition.

\subsection{Rupture-size laws}

The sum of $k$ i.i.d.\ geometric variables is negative binomial, so
\eqref{eq:rkdecomp} gives immediately:
\begin{equation}
\pr{r_k=n}=\mathcal{A}^{(n)}_{k}\,s^{k}\,\rho^{\,n-1},
\qquad
\mathcal{A}^{(n)}_{k}=\binom{n+k-2}{k-1},
\quad n\ge1.
\label{eq:rklaw}
\end{equation}
The first few cases display the pattern:
\begin{equation}
\pr{r_1=n}=s\rho^{n-1}\ \text{(geometric --- matches one-cell BDC)},
\end{equation}
\begin{equation}
\pr{r_2=n}=n\,s^2\rho^{n-1},
\qquad
\pr{r_3=n}=\tfrac{n(n+1)}{2}s^3\rho^{n-1},
\qquad
\pr{r_4=n}=\tfrac{n(n+1)(n+2)}{6}s^4\rho^{n-1}.
\end{equation}

Equivalently, the coefficients satisfy the recurrence
\begin{equation}
\mathcal{A}^{(n)}_{k}=\sum_{i=1}^{n}\mathcal{A}^{(i)}_{k-1},
\qquad
\mathcal{A}^{(i)}_{1}=1,
\end{equation}
which is just \eqref{eq:rklaw} conditioned on the value of $r_{k-1}$:
$\pr{r_k=n}=\sum_i\pr{r_{k-1}=i}\pr{G_k=n-i}$. The closed form
$\mathcal{A}^{(n)}_{k}=\binom{n+k-2}{k-1}$ then follows by induction from
the hockey-stick identity --- supplying the proof that the original
investigation had conjectured but not completed.

\begin{remark}[Correction of the draft's closed form]
The draft recorded $\mathcal{A}^{(n)}_{k}=\frac{1}{k!}\frac{(n+k-1)!}{(n-1)!}
=\binom{n+k-1}{k}$, a formula off by one shift: it contradicts the draft's
own $r_2$ and $r_3$ cases, which \eqref{eq:rklaw} reproduces exactly. The
correct coefficient is $\binom{n+k-2}{k-1}$.
\end{remark}

Further consequences, all verified numerically (see
\S\ref{sec:chained-verification}):
\begin{align}
\mean{r_k}&=1+k\,\frac{\rho}{s}=1+k\,\frac{\beta}{\delta},
&
\mathrm{Var}(r_k)&=k\,\frac{\rho}{s^2},
\label{eq:rkmoments}\\[0.4em]
\mean{z^{r_k}}&=z\Bigl(\frac{s}{1-\rho z}\Bigr)^k,
&
\sum_{n\ge1}\pr{r_k=n}&=1\ \text{(burst is certain for $\mu=0$)}.
\label{eq:rkgf}
\end{align}
The mean rupture size grows linearly along the chain --- each cell adds one
expected burst's worth of births, $\beta/\delta$, to whatever it inherited.

\subsection{Rupture times}
\label{sec:rupture-times}

The first rupture time inherits the single-cell flux. With $m$ founders, the
no-rupture probability is $I(t)^m$ by \eqref{eq:Ihatk} (with $\mu=0$,
$D\equiv0$), so
\begin{equation}
\pr{t_1>t}=I(t)^m,
\qquad
f_{t_1}^{(m)}(t)=-m\,I(t)^{m-1}I'(t)=m\,I(t)^{m-1}\delta J(t).
\label{eq:t1m}
\end{equation}
For $m=1$ this is the burst-time density $\varphi=\delta J$ of
\S\ref{sec:burst-time} (which is proper here, since $b=0$). When $\mu>0$,
the same reasoning distinguishes the two fixation routes, as in the
original investigation: the \emph{burst-only} time has CDF
$(1-I(t))/(1-b)$ and density $\delta J(t)/(1-b)$, while the
\emph{general-fixation} time --- rupture or internal extinction --- has
density
\begin{equation}
\delta J(t)+\mu\,p_1(t),
\end{equation}
the extra term being the flux through $1\to0$: only a population of size
exactly one can die out in a single death event.

\subsection{Inter-rupture intervals}

Given $k$ founders, the interval $T(k)$ until rupture decomposes over the
holding times: if $G=g$ births precede the catastrophe, the load runs
through $k,k+1,\ldots,k+g$, so
\begin{equation}
T(k)=\sum_{j=k}^{k+G}\xi_j,
\qquad
\xi_j\sim\mathrm{Exp}\lrb{(\beta+\delta)j}\ \text{independent}.
\end{equation}
Averaging over $G\sim\mathrm{Geom}_0(s)$, the Laplace transform is
\begin{equation}
\Lap{T}(k,u)=\mean{e^{-uT(k)}}
=\sum_{g\ge0}s\rho^g\prod_{j=k}^{k+g}\frac{(\beta+\delta)j}{(\beta+\delta)j+u}.
\label{eq:LTk}
\end{equation}
Writing $u'=u/(\beta+\delta)$ and unrolling the product with Gamma
functions, the series sums to a hypergeometric form, continuing the theme of
\S\ref{sec:peff-section}:
\begin{equation}
\Lap{T}(k,u)=\frac{s\,k}{k+u'}\,\Fhyp\Bigl(k+1,1;k+1+u';\rho\Bigr).
\label{eq:LTk-hyp}
\end{equation}
Differentiating \eqref{eq:LTk} at $u=0$ gives the exact mean
\begin{equation}
\mean{T(k)}=\sum_{m\ge0}\frac{\rho^m}{(\beta+\delta)(k+m)},
\label{eq:meanT}
\end{equation}
which is strictly decreasing in $k$: the more a cell inherits, the faster it
bursts. This quantifies the original investigation's qualitative
observation that the chain is ``partitioned by reducing intervals of time''.
For $(\beta,\delta)=(1,1)$: $\mean{T(1..4)}\approx0.693,0.386,0.273,0.212$
for fixed founding loads; the chain intervals $T_k$, whose founding loads
are random, decrease faster still
($0.691,0.500,0.377,0.292$ in simulation).

\subsection{The second rupture time}

The chain's dependence structure is now visible. The second rupture time is
$t_2=t_1+T_2$, but $t_1$ and $T_2$ are \emph{not} independent: both are
functions of $G_1$, the number of births in the first cell. Conditional on
$G_1=g$, however, the two pieces involve disjoint sets of exponential
holding times, so their Laplace transforms factor:
\begin{equation}
\mean{e^{-ut_2}\mid G_1=g}
=P_1(1+g,u)\,\Lap{T}(1+g,u),
\qquad
P_1(m,u)=\prod_{j=1}^{m}\frac{(\beta+\delta)j}{(\beta+\delta)j+u}.
\end{equation}
Averaging over $G_1\sim\mathrm{Geom}_0(s)$:
\begin{equation}
\Lap{t_2}(u)
=\sum_{g\ge0} s\rho^g\,P_1(1+g,u)\,\Lap{T}(1+g,u).
\label{eq:Lt2}
\end{equation}
This resolves the question with which the original investigation closed. The
convolution guessed there,
$\rho(t_2)=\int f(\tau)\sum_k\pr{r_1=k|\tau}\phi_k(t_2-\tau)\dd\tau$,
is the same statement in different clothes --- and, as suspected, it is not
the tractable one. The Laplace form \eqref{eq:Lt2} is; so is the density
itself, computable exactly as a mixture of hypoexponential densities
($t_1|G_1$ and $T_2|(G_1,G_2)$ are sums of independent exponentials) and
verified against simulation in \S\ref{sec:chained-verification}. Moments
assemble from the pieces, e.g.\
$\mean{t_2}=\mean{T(1)}+\sum_{n\ge1}\pr{r_1=n}\mean{T(n)}$.

\subsection{Verification record and scope}
\label{sec:chained-verification}

Every quantitative claim of this section was checked by exact simulation of
the chained process (the event-type decomposition makes sampling direct):
the suite \texttt{verify\_chained\_transfer.py} comprises 28 checks across
three parameter sets --- the coefficient identity (exact integer
arithmetic); the rupture-size laws for $k=1..4$ against $10^6$ chains;
moments \eqref{eq:rkmoments}; the \pgf{} \eqref{eq:rkgf}; the $t_1$ density
\eqref{eq:t1m} vs $\delta J$; the interval transform \eqref{eq:LTk} for
fixed founders $k\in\{1,2,3,5\}$; the random-founder mixture consistency of
the chain intervals; the $t_2$ transform \eqref{eq:Lt2}; the exact
hypoexponential-mixture density of $t_2$ (relative $L^2\approx3\%$ against
histograms); and the monotone interval means. \textbf{Final run 2026-08-07:
28/28 PASS.}

What remains open, and is stated plainly: the joint laws of $(t_k,r_k)$ for
general $\mu>0$ (the $\Theta/\Omega$ coefficient machinery of the draft,
with its ${}_2F_1$ antiderivatives, is the route; see
Appendix~\ref{app:technical}), and the joint law of $(t_M,r_M)$ for the full
$M$-cell chain. For $\mu=0$, the decomposition \eqref{eq:rkdecomp} extends
the size laws to all $M$ at once, since nothing about it depends on where a
cell sits in the chain.
